%%%%%%%%%%%%%%%%%%%%%%%%%%%%%%%%%%%%%%%%%%%%%%%%%%%%%%%%%%%%%%%%%%%%%%%%%%%%%%%%%%
\begin{frame}[fragile]\frametitle{Notes of Simon Dixon}
\begin{center}
{\Large Geopolitics}
\end{center}
\end{frame}
%%%%%%%%%%%%%%%%%%%%%%%%%%%%%%%%%%%%%%%%%%%%%%%%%%%%%%%%%%%%
\begin{frame}[fragile]\frametitle{Why This Talk Matters}
  \begin{itemize}
    \item Every major geopolitical event has two layers: the visible narrative and the financial logic
    \item Headlines optimise for drama; markets optimise for profit
    \item The gap between the two is where the real signal lives
    \item The Iran conflict is a live case study in reading that gap
    \item Dixon's framework: follow the money, not the missiles
    \item Aim: give you a durable mental model, not just a news summary
    \item Held cautiously, a map, not a photograph
  \end{itemize}
\end{frame}

%%%%%%%%%%%%%%%%%%%%%%%%%%%%%%%%%%%%%%%%%%%%%%%%%%%%%%%%%%%%
\begin{frame}[fragile]\frametitle{Key Terms: The Economic Vocabulary}
  \begin{itemize}
    \item \textbf{K-shaped economy}: recovery splits: asset owners rise, wage earners fall
    \item \textbf{Petrodollar}: 1970s deal: oil priced in USD creates permanent dollar demand
    \item \textbf{Financialisation}: financial motives dominate over production and trade
    \item \textbf{Fiscal dominance}: debt so large that monetary policy serves fiscal needs
    \item \textbf{De-dollarisation}: countries reducing USD dependence in trade and reserves
    \item \textbf{Force majeure}: contract clause suspending obligations under extraordinary events
    \item \textbf{Strategic Petroleum Reserve (SPR)}: government emergency oil stockpiles
    \item \textbf{Private credit}: non-bank lending; illiquid, opaque, growing fast
  \end{itemize}
\end{frame}

%%%%%%%%%%%%%%%%%%%%%%%%%%%%%%%%%%%%%%%%%%%%%%%%%%%%%%%%%%%%
\begin{frame}[fragile]\frametitle{Key Terms: The Power Structures}
  \begin{itemize}
    \item \textbf{MIC}: Military-Industrial Complex; profits from permanent conflict
    \item \textbf{FIC}: Financial-Industrial Complex; profits from financing and rebuilding
    \item \textbf{TIC}: Technical-Industrial Complex; profits from data, AI, surveillance
    \item \textbf{Multipolar world}: power distributed across US, China, Russia, Gulf, India
    \item \textbf{Proof of Weapons Network}: Dixon's term for FIC + MIC + TIC as interlocking system
    \item \textbf{Self-custody}: holding your own Bitcoin keys; financial sovereignty in practice
    \item \textbf{Multipolar transition}: managed or contested shift away from US unipolarity
  \end{itemize}
\end{frame}

%%%%%%%%%%%%%%%%%%%%%%%%%%%%%%%%%%%%%%%%%%%%%%%%%%%%%%%%%%%%
\begin{frame}[fragile]\frametitle{500 Years: The Imperial Playbook}
  \begin{itemize}
    \item \textbf{1600s:} East India Companies: private, state-backed, militarised multinationals
    \item Privatised profits of empire; socialised military costs through state power
    \item Template: extract from periphery, concentrate wealth at centre
    \item \textbf{1800s:} British Empire: sterling as reserve currency, Bank of England at hub
    \item Naval supremacy enforced the financial arrangement
    \item When empire became too costly, same networks profited from contraction
    \item Financed the wars ending British hegemony; positioned for American dominance
    \item \alert{Pattern: financial institutions survive every hegemonic transition}
  \end{itemize}
\end{frame}

%%%%%%%%%%%%%%%%%%%%%%%%%%%%%%%%%%%%%%%%%%%%%%%%%%%%%%%%%%%%
\begin{frame}[fragile]\frametitle{From Gold Standard to Petrodollar}

% \begin{tikzpicture}[x=1.95cm, y=1cm, font=\small]
  % % Timeline axis
  % \draw[->, thick, color=darkblue] (0,0) -- (6.5,0);

  % % Nodes
  % \foreach \x/\yr/\lbl in {
    % 0.3/1913/{Fed Reserve\\created},
    % 1.3/1944/{Bretton Woods\\dollar=gold anchor},
    % 2.3/1971/{Nixon closes\\gold window},
    % 3.3/1973/{Petrodollar\\deal (Saudi)},
    % 4.3/2001/{China WTO\\manufacturing shift},
    % 5.3/2008/{Financial crisis\\bailout era},
    % 6.3/Now/{Unipolar\\era ending}
  % }{
    % \draw[accent, thick] (\x,0) -- (\x,0.35);
    % \node[above, align=center, text=darkblue, font=\tiny] at (\x,0.35) {\textbf{\yr}\\\lbl};
  % }
% \end{tikzpicture}

  % \vspace{0.4em}
  \begin{itemize}
    \item 1971 Nixon shock: dollar untethered from gold: effective sovereign default
    \item 1973--74 petrodollar: oil in USD only $\Rightarrow$ permanent global dollar demand
    \item Freed the US from any external monetary constraint
    \item Manufacturing offshored; capital markets grew; debt became growth engine
  \end{itemize}
\end{frame}

%%%%%%%%%%%%%%%%%%%%%%%%%%%%%%%%%%%%%%%%%%%%%%%%%%%%%%%%%%%%
\begin{frame}[fragile]\frametitle{The Hollowing Out: 1990s to Present}
  \begin{itemize}
    \item China's WTO entry (2001) accelerated offshoring of Western manufacturing at scale
    \item China absorbed the industrial base; West kept the financial intermediation fees
    \item 2008: failure of mortgage-backed securities nearly collapsed the global system
    \item Bailout transferred enormous wealth from taxpayers to financial institutions
    \item China GDP: from 30\% of US to approaching parity
    \item US share of global GDP: from over 40\% to around 25\%
    \item Gulf SWFs shifted from buying US Treasuries to buying US equities
    \item BRICS expanded; dollar's share of global reserves declined
  \end{itemize}
\end{frame}

%%%%%%%%%%%%%%%%%%%%%%%%%%%%%%%%%%%%%%%%%%%%%%%%%%%%%%%%%%%%
\begin{frame}[fragile]\frametitle{The Three Power Complexes}
  \begin{itemize}
    \item \textbf{MIC} (Military-Industrial Complex)
      \begin{itemize}
        \item Economic model requires ongoing threat, spending, and war
        \item War is not a failure state: war is the product
        \item Wants: forever war
      \end{itemize}
    \item \textbf{FIC} (Financial-Industrial Complex)
      \begin{itemize}
        \item Profits from financing wars, managing capital flows, reconstruction
        \item Forever war is bad for asset prices
        \item Wants: managed transition to stable, profitable new order
      \end{itemize}
    \item \textbf{TIC} (Technical-Industrial Complex)
      \begin{itemize}
        \item Profits from data, attention, and behavioural control
        \item Provides the surveillance layer atop MIC and FIC
        \item Wants: architecture of a managed population
      \end{itemize}
  \end{itemize}
\end{frame}

%%%%%%%%%%%%%%%%%%%%%%%%%%%%%%%%%%%%%%%%%%%%%%%%%%%%%%%%%%%%
\begin{frame}[fragile]\frametitle{The FIC Is Currently Winning}
  \begin{itemize}
    \item Dixon's core thesis: Iran conflict is a \emph{controlled demolition} of the old order
    \item Rubble already priced into contracts before the dust settles
    \item FIC does not need to fire a missile: it pulls insurance coverage
    \item Nation states function as proxies; the real war is between power factions
    \item Gulf sovereign wealth funds now co-investors in the outcome, not subordinates
    \item Saudi Aramco holds a board seat at BlackRock
    \item Jared Kushner's Affinity Partners backed by Gulf sovereign wealth funds (public record)
    \item The transition is managed: messy, contested, but directionally deliberate
  \end{itemize}
\end{frame}

%%%%%%%%%%%%%%%%%%%%%%%%%%%%%%%%%%%%%%%%%%%%%%%%%%%%%%%%%%%%
\begin{frame}[fragile]\frametitle{Market Signals Don't Lie}
  \begin{center}
  \begin{tabular}{llll}
    \toprule
    \textbf{Asset} & \textbf{WW3 Expected} & \textbf{Reality} & \textbf{Reading} \\
    \midrule
    Gold       & $>$\$6,000/oz      & Barely moved         & No panic \\
    Oil        & $>$\$150, no ceil  & \$115 $\to$ \$89     & Managed cap \\
    S\&P 500   & Severe decline     & Modest pullback      & No collapse \\
    Dollar     & Sharp weakening    & Held                 & No flight \\
    \bottomrule
  \end{tabular}
  \end{center}
  \vspace{0.5em}
  \begin{itemize}
    \item Sophisticated money is not positioned for Armageddon
    \item It is positioned for an off-ramp
    \item Markets pricing a deal, not a world war
    \item Gold is the canonical flight-to-safety asset: its flatness is a signal
  \end{itemize}
\end{frame}

%%%%%%%%%%%%%%%%%%%%%%%%%%%%%%%%%%%%%%%%%%%%%%%%%%%%%%%%%%%%
\begin{frame}[fragile]\frametitle{Hormuz: Closed by Insurance, Not Missiles}
  \begin{itemize}
    \item 20\% of global oil and significant LNG passes through the Strait of Hormuz
    \item Strait not physically blockaded by any naval force
    \item War risk underwriters withdrew marine coverage: ships stopped commercially
    \item Aviation insurance was \emph{not} pulled: planes kept flying through same airspace
    \item Ships and planes face materially different operational risk profiles (counterargument)
    \item But: asymmetry implies deliberate choice, not pure commercial logic
    \item Dixon: the financial system did what missiles could not
    \item Single most compelling piece of evidence in the entire framework
  \end{itemize}
\end{frame}

%%%%%%%%%%%%%%%%%%%%%%%%%%%%%%%%%%%%%%%%%%%%%%%%%%%%%%%%%%%%
\begin{frame}[fragile]\frametitle{The Russia Anomaly}
  \begin{itemize}
    \item Oil spiked from mid-\$50s to \$115 on conflict outbreak
    \item Crashed back to \$89 within days
    \item Three drivers of the reversal: Trump announcements, G7 SPR pledges, Russian oil
    \item US Treasury \alert{quietly lifted sanctions on Russian oil} to suppress the price
    \item America went to war with Iran and relieved Russia simultaneously
    \item The price ceiling was the priority: not the geopolitical narrative
    \item Counterargument: pragmatic crisis management, not pre-arranged script
    \item But: decision made at speed, at the highest Treasury level: hard to dismiss
  \end{itemize}
\end{frame}

%%%%%%%%%%%%%%%%%%%%%%%%%%%%%%%%%%%%%%%%%%%%%%%%%%%%%%%%%%%%
\begin{frame}[fragile]\frametitle{Russia: Clearest Short-Term Winner}
  \begin{itemize}
    \item Pre-war Russian oil: $\approx$\$35--40/barrel (sanctions discount)
    \item During conflict: approaching \$100/barrel
    \item Near-tripling of revenue with zero military involvement and zero additional cost
    \item Windfall recycled into gold, commodities, and BRICS partnerships: not US Treasuries
    \item Russia as swing oil supplier: direct leverage over European energy security
    \item Europe chose: Qatari LNG (now constrained) and US LNG (on predatory terms)
    \item Alternative: renegotiate with Russia: either way, leverage shifts away from the West
    \item Caveat: Russia still losing demographic and technological competition long-term
  \end{itemize}
\end{frame}

%%%%%%%%%%%%%%%%%%%%%%%%%%%%%%%%%%%%%%%%%%%%%%%%%%%%%%%%%%%%
\begin{frame}[fragile]\frametitle{The Gulf States: From Creditors to Co-Owners}
  \begin{itemize}
    \item Classic role: buy US Treasury bonds (creditor, subordinate)
    \item New role: buy US equities, take board seats, co-invest (co-owner, stakeholder)
    \item Saudi Aramco board seat at BlackRock (public record)
    \item Gulf capital now funding AI infrastructure and global asset management
    \item Structural shift documented and verifiable: not speculative
    \item What the Gulf wants from this leverage:
      \begin{itemize}
        \item Regional stability
        \item Exit of US military bases
        \item Normalisation of Iran
        \item Resolution of the Palestinian cause
      \end{itemize}
    \item They are not seeking more war: war is bad for their investment portfolios
  \end{itemize}
\end{frame}

%%%%%%%%%%%%%%%%%%%%%%%%%%%%%%%%%%%%%%%%%%%%%%%%%%%%%%%%%%%%
\begin{frame}[fragile]\frametitle{China Made Itself Indispensable}
  \begin{itemize}
    \item Became the manufacturing base for everyone, including rivals
    \item US F-35 jets require Chinese rare earth minerals
    \item American munitions replenishment requires Chinese supply chains
    \item European energy infrastructure requires Chinese components
    \item China brokered Iran--Saudi normalisation in 2023: a genuine diplomatic earthquake
    \item China now outproduces US in: shipbuilding, EVs, solar panels, advanced electronics
    \item US share of global GDP: $>$40\% (peak) $\to$ $\approx$25\% today
    \item China GDP: 30\% of US (2008) $\to$ approaching parity
  \end{itemize}
\end{frame}

%%%%%%%%%%%%%%%%%%%%%%%%%%%%%%%%%%%%%%%%%%%%%%%%%%%%%%%%%%%%
\begin{frame}[fragile]\frametitle{The Multipolar Transition}
  \begin{itemize}
    \item Unipolar era (post-1991): US dominated military, financial, and reserve currency systems
    \item That system required the Middle East in perpetual tension: unified Gulf was a threat
    \item Emerging order: US, China, Russia, India, Gulf as distinct power centres
    \item China's rise changed the equation irreversibly
    \item Raoul Pal, Jordi Visser, and mainstream macro analysts reach the same structural conclusion
    \item Direction not seriously in doubt; pace and coordination are the open questions
    \item Transitions of this scale are messy, contested, full of unintended consequences
    \item The destination may be correct even if the route map is wrong
  \end{itemize}
\end{frame}

%%%%%%%%%%%%%%%%%%%%%%%%%%%%%%%%%%%%%%%%%%%%%%%%%%%%%%%%%%%%
\begin{frame}[fragile]\frametitle{Private Credit: The Hidden Stress Signal}
  \begin{itemize}
    \item Retail investors sold: institutional-grade returns with quarterly liquidity windows
    \item Force majeure clauses buried in the small print
    \item War declared $\Rightarrow$ clauses invoked $\Rightarrow$ money locked
    \item Funds unable to meet redemptions selling assets at steep discounts
    \item Some bids as low as \alert{14 cents on the dollar} (Dixon, cited claim)
    \item Underlying assets: companies being disrupted by AI
    \item SaaS cashflows eroding as AI automates core functions
    \item Double squeeze: deteriorating assets + rising redemption pressure simultaneously
  \end{itemize}
\end{frame}

%%%%%%%%%%%%%%%%%%%%%%%%%%%%%%%%%%%%%%%%%%%%%%%%%%%%%%%%%%%%
\begin{frame}[fragile]\frametitle{Abundance vs Scarcity: The Portfolio Imbalance}
  \begin{itemize}
    \item S\&P 500: $\approx$50\% technology and communications
    \item S\&P 500: only $\approx$5\% materials and energy combined
    \item Years of capital flowing towards \emph{abundance} assets over \emph{scarcity} assets
    \item AI rapidly commoditises software: abundance accelerates
    \item Physical inputs to the intelligence economy become the scarce premium
    \item Rebalancing when it comes will be painful for those on the wrong side
    \item This is not a cyclical trade: it is a structural decade-long rotation
    \item Pal and Visser analysis supports the same directional conclusion
  \end{itemize}
\end{frame}

%%%%%%%%%%%%%%%%%%%%%%%%%%%%%%%%%%%%%%%%%%%%%%%%%%%%%%%%%%%%
\begin{frame}[fragile]\frametitle{Where to Position: The Scarcity Thesis}
  \begin{itemize}
    \item \textbf{Energy infrastructure}: integrated energy, midstream, LNG terminals; essential in contested supply chains
    \item \textbf{Critical minerals}: copper, silver, rare earths; China dominates refining; decade-long squeeze
    \item \textbf{European defence}: most durable budget cycle in a generation; energy security broadens the thesis
    \item \textbf{Gold}: not a trade; hedge against the endgame; Global South central banks accumulating at record rates
    \item \textbf{Bitcoin (self-custody)}: force majeure on private credit is a preview of capital controls
    \item Common thread across all five: \alert{physical scarcity}
    \item Question is not whether to position in scarcity: it is how much, and how patiently
  \end{itemize}
\end{frame}

%%%%%%%%%%%%%%%%%%%%%%%%%%%%%%%%%%%%%%%%%%%%%%%%%%%%%%%%%%%%
\begin{frame}[fragile]\frametitle{Holding the Framework Honestly}
  \begin{itemize}
    \item \textbf{Strongest evidence:} insurance asymmetry: ships stopped, planes flew
    \item \textbf{Strong evidence:} Russia sanctions relief at speed, at highest Treasury level
    \item \textbf{Strong evidence:} market behaviour inconsistent with genuine WW3 pricing
    \item \textbf{Strong evidence:} Gulf capital shift from Treasuries to equities (public, verifiable)
    \item \textbf{Strong evidence:} China's supply chain indispensability (factual, documented)
    \item \textbf{Strong evidence:} Iran--Saudi normalisation brokered by China in 2023
    \item Dixon's conflict is moving faster and more unpredictably than he anticipated (his own admission)
    \item Treat as structural orientation: not a verified prediction
  \end{itemize}
\end{frame}

%%%%%%%%%%%%%%%%%%%%%%%%%%%%%%%%%%%%%%%%%%%%%%%%%%%%%%%%%%%%
\begin{frame}[fragile]\frametitle{Where to Challenge Dixon}
  \begin{itemize}
    \item MIC vs FIC risks being unfalsifiable: any outcome can be retrofitted post-hoc
    \item Coordinated intent across thousands of independent actors strains credibility
    \item Aviation and marine insurance are genuinely separate markets with different risk models
    \item Russia sanctions relief may be pragmatic crisis management, not pre-arranged
    \item Gulf states are rival monarchies with conflicting interests: not a unified bloc
    \item China has structural vulnerabilities: ageing population, property crisis, political concentration
    \item Private credit stress predates this conflict; mundane causes exist (overleverage, duration mismatch)
    \item Markets can be wrong, manipulated, or slow to price tail risks (2008, COVID precedents)
  \end{itemize}
\end{frame}

%%%%%%%%%%%%%%%%%%%%%%%%%%%%%%%%%%%%%%%%%%%%%%%%%%%%%%%%%%%%
\begin{frame}[fragile]\frametitle{Three Questions for Any Geopolitical Event}
  \begin{itemize}
    \item \textbf{Who profits from the chaos?}
      \begin{itemize}
        \item Follow the contracts, capital flows, insurance adjustments: not the stated motives
      \end{itemize}
    \item \textbf{What are the markets actually pricing?}
      \begin{itemize}
        \item Gold, oil, equities, and credit spreads often know things the headlines do not
      \end{itemize}
    \item \textbf{What is the off-ramp?}
      \begin{itemize}
        \item Every managed conflict needs one: who is positioned to offer it, and at what price?
      \end{itemize}
    \item \textbf{Bonus:} read the insurance market before reading the newspaper
    \item \textbf{Bonus:} know your force majeure clauses before a crisis, not during one
    \item The unipolar dollar era is ending: position for a multipolar world
    \item The fog of financial misdirection is thicker than the fog of war
  \end{itemize}
\end{frame}

%%%%%%%%%%%%%%%%%%%%%%%%%%%%%%%%%%%%%%%%%%%%%%%%%%%%%%%%%%%%
\begin{frame}[fragile]\frametitle{Closing}
  \begin{center}
    \large
    \textit{``The fog of war is real.}\\
    \textit{The fog of financial misdirection is thicker still.''}\\[1em]
    \normalsize
    Following the money does not guarantee you will always be right\\[0.3em]
    But it is a far better compass than following the headlines\\[1.5em]
    \textbf{accent} Framework: Simon Dixon $|$ Hard Talk Live, March 2026\\
    Additional macro context: Raoul Pal $\cdot$ Jordi Visser
  \end{center}
\end{frame}
